\headline{{\textbf{\Large\sffamily Seasonal Care Tips:}}\\
          {\textbf{\large\sffamily Mid\--June to Mid\--July}}}
\label{2025-06-care}
\noindent
From mid\--June to mid\--August, bonsai trees in the London area experience the height of summer.
With long daylight hours and typically warm temperatures, growth remains vigorous, but the season also brings challenges---intense sun, potential drought stress, and persistent pests.
Balancing growth refinement with protection from the elements is key.
Here's how to guide your trees through these peak summer months.

\medskip
\topic{Refining and Controlling Growth}
\noindent
Deciduous species like maples, beeches, and elms will now have dense foliage.
Continue selective pruning to maintain shape and encourage finer branching, but avoid heavy defoliation during heatwaves, as this can stress the tree.
Instead, thin out overcrowded areas to improve airflow and light penetration.

Pines that were decandled earlier should now produce their second flush.
Monitor new shoots and adjust wiring if needed.
Junipers and other conifers benefit from periodic trimming to maintain foliage pads---avoid cutting back to bare wood, as summer heat can hinder recovery.
For flowering species like azaleas, remove spent blooms to redirect energy into healthy growth.

\medskip
\topic{Watering: The Critical Routine}
\noindent
Summer demands meticulous watering.
Check soil moisture twice daily---morning and evening---especially for small pots or fast\--draining mixes.
Water deeply until it runs freely from the drainage holes, ensuring the entire root mass is hydrated.
Avoid shallow sprinkling, which encourages surface roots and weakens the tree.

During heatwaves, consider grouping trees together or placing pots on damp gravel trays to increase humidity.
For those away for short periods, capillary matting or self\--watering systems can help, but test them beforehand to avoid over\-- or under\--watering.

\medskip
\topic{Fertilising for Sustained Health}
\noindent
Reduce nitrogen\--heavy fertilisers for deciduous trees by late July to harden off new growth before autumn.
Conifers still benefit from balanced feeding to strengthen needles and roots.
Slow\--release pellets or liquid feeds diluted to half\--strength are ideal in high temperatures to prevent root burn.

For flowering or fruiting bonsai, switch to a phosphorus\--rich fertiliser to support bud formation for next year.
Organic options like fish emulsion or compost tea can be used sparingly to avoid attracting pests.
Always water thoroughly before applying fertiliser to protect roots.

\columnbreak

\medskip
\topic{Pestling with Pests and Diseases}
\noindent
Aphids, spider mites, and scale insects thrive in warm weather.
Inspect foliage regularly for discolouration, sticky residue, or fine webbing.
Blast pests off with water or treat with insecticidal soap, applying in the early morning or evening to avoid leaf scorch.

Fungal infections like powdery mildew can arise in humid conditions.
Improve airflow by spacing trees and pruning dense growth.
Remove affected leaves promptly and apply a fungicide if necessary.
Avoid overhead watering in the evening to keep foliage dry overnight.

\medskip
\topic{Sun Protection and Positioning}
\noindent
While most hardy species tolerate full sun, delicate trees like Japanese maples or young saplings may need shade during peak afternoon heat.
Use shade cloth or position them under a pergola.
Tropical species (ficus, serissa, etc.) relish the warmth but should be sheltered from strong winds that can dry out foliage.

Monitor pot temperatures---light\--coloured pots or stands can help reflect heat.
Elevate pots slightly to improve airflow and prevent overheating from paving stones.

\medskip
\topic{Wiring and Structural Adjustments}
\noindent
Summer growth can cause wires to bite into bark quickly.
Inspect wired branches weekly and loosen if necessary.
For deciduous trees, consider removing wires after 2\--3 months to avoid scarring.
Conifers can often retain wiring longer but check for swelling.

Avoid major styling or repotting now---save heavy work for cooler months.
Focus instead on minor tweaks and maintenance to keep trees healthy and balanced.

\medskip
\topic{Preparing for Late Summer and Autumn}
\noindent
By mid\--August, start planning for autumn.
Gradually reduce fertilisation to let trees prepare for dormancy.
Begin collecting seeds or cuttings for propagation, and clean tools and benches to prevent disease carryover.

If moss has dried out, rehydrate it gently or replace it to maintain a tidy surface.
Mulch can be added to retain moisture but keep it away from the trunk to avoid rot.

\medskip
\topic{Enjoying the Summer Bonsai Life}
\noindent
Take time to appreciate the lush vitality of your trees.
Document their progress with photos, noting how they've responded to your care.
Share observations with the club---whether it's a perfectly ramified branch or a lesson learned from a pest battle.
Summer is a season of abundance and challenge, but with attentive care, your bonsai will thrive.
\closearticle
\pagebreak
